% Adapted from Cyril Langlois; CC BY-SA 4.0. See ATTRIBUTION.md.
\documentclass[tikz,border=8pt]{standalone}
\usepackage{minimalist-profile}
\begin{document}
\begin{tikzpicture}[ms base,x={(0.866cm,-0.5cm)}, y={(0.866cm,0.5cm)}, z={(0cm,1cm)}, scale=1.0,
    %Option for nice arrows
    % Native Latex heads come from the shared profile.
    inner sep=\msSpaceLabel, outer sep=0bp,%
    axis/.style={ms arrow,draw=msMuted,line width=\msLineContext},
    wave/.style={ms structure,color=#1,smooth},
    polaroid/.style={draw=msMuted,line width=\msLineContext,fill=msMuted,ms fill},
    field vector/.style={line width=\msLineContext,-{Latex[length=0.6*\msArrowLength,width=0.6*\msArrowWidth]}},
]
    % Shared colors are reused for waves and labels.
    % Frame
    \coordinate (O) at (0, 0, 0);
    \draw[axis] (O) -- +(14, 0,   0) node [right,text=msMuted] {$x$};
    \draw[axis] (O) -- +(0,  2.5, 0) node [right,text=msMuted] {$y$};
    \draw[axis] (O) -- +(0,  0,   2) node [above,text=msMuted] {$z$};

    \draw[ms structure,ms dashed] (-2,0,0) -- (O);

    % monochromatic incident light with electric field
    \draw[wave=msColor6, draw opacity=\msOpacityForeground, variable=\x, samples at={-2,-1.75,...,0}]
        plot (\x, { cos(1.0*\x r)*sin(2.0*\x r)}, { sin(1.0*\x r)*sin(2.0*\x r)})
        plot (\x, {-cos(1.0*\x r)*sin(2.0*\x r)}, {-sin(1.0*\x r)*sin(2.0*\x r)});

    \foreach \x in{-2,-1.75,...,-0.25}{
        \draw[color=msColor6, draw opacity=\msOpacityForeground,field vector]
            (\x,0,0) -- (\x, { cos(1.0*\x r)*sin(2.0*\x r)}, { sin(1.0*\x r)*sin(2.0*\x r)})
            (\x,0,0) -- (\x, {-cos(1.0*\x r)*sin(2.0*\x r)}, {-sin(1.0*\x r)*sin(2.0*\x r)});
    }

    \filldraw[polaroid] (0,-2,-1.5) -- (0,-2,1.5) -- (0,2,1.5) -- (0,2,-1.5) -- (0,-2,-1.5)
        ;
    \node[anchor=north west] at (0,-2,-1.7) {Polaroid};

    %Direction of polarization
    \draw[ms structure,<->] (0,-1.75,-1) -- (0,-0.75,-1);

    % Electric field vectors
    \draw[wave=msColor6, variable=\x,samples at={0,0.25,...,6}]
        plot (\x,{sin(2*\x r)},0)node[anchor=north,text=msColor6]{$\vec{E}$};

    %Polarized light between polaroid and thin section
    \foreach \x in{0.25,0.5,...,6}
        \draw[color=msColor6,field vector] (\x,0,0) -- (\x,{sin(2*\x r)},0);

    \draw (3,1,1) node [text width=2.5cm, text centered,text=msColor6]{Polarized light};

    %Crystal thin section
    \begin{scope}[draw=msMuted,line width=\msLineContext]
        \draw (6,-2,-1.5) -- (6,-2,1.5)
                -- (6, 2, 1.5) -- (6, 2, -1.5) -- cycle % First face
            (6,  -2, -1.5) -- (6.2, -2,-1.5)
            (6,   2, -1.5) -- (6.2,  2,-1.5)
            (6,  -2,  1.5) -- (6.2, -2, 1.5)
            (6,   2,  1.5) -- (6.2,  2, 1.5)
            (6.2,-2, -1.5) -- (6.2, -2, 1.5) -- (6.2, 2, 1.5)
                -- (6.2, 2, -1.5) -- cycle; % Second face

        \node[anchor=north west] at (6,-2,-1.7) {Crystal section};

        %Optical indices
        \draw[msColor1, ->]       (6.1, 0, 0) -- (6.1, 0.26,  0.966) node [right,text=msColor1] {$n_{g}'$}; % index 1
        \draw[msColor1, ms dashed]   (6.1, 0, 0) -- (6.1,-0.26, -0.966); % index 1
        \draw[msColor5, ->]     (6.1, 0, 0) -- (6.1, 0.644,-0.173) node [right,text=msColor5] {$n_{p}'$}; % index 2
        \draw[msColor5, ms dashed] (6.1, 0, 0) -- (6.1,-0.644, 0.173); % index 2
    \end{scope}

    %Rays leaving thin section
    \draw[wave=msColor1,   variable=\x, samples at={6.2,6.45,...,12}]
        plot (\x, {0.26*0.26*sin(2*(\x-0.5) r)},  {0.966*0.26*sin(2*(\x-0.5) r)});  %n'g-oriented ray
    \draw[wave=msColor5, variable=\x, samples at={6.2,6.45,...,12}]
        plot (\x, {0.966*0.966*sin(2*(\x-0.1) r)},{-0.26*0.966*sin(2*(\x-0.1) r)}); %n'p-oriented ray
    \draw (10,1,1) node [text width=2.5cm, text centered] {Polarized and dephased light};

    \foreach \x in{6.2,6.45,...,12} {
        \draw[color=msColor5,field vector] (\x, 0, 0) --
            (\x, {0.966*0.966*sin(2*(\x-0.1) r)}, {-0.26*0.966*sin(2*(\x-0.1) r)});
        \draw[color=msColor1,field vector] (\x, 0, 0) --
            (\x, {0.26*0.26*sin(2*(\x-0.5) r)}, {0.966*0.26*sin(2*(\x-0.5) r)});
    }

    %Second polarization
    \draw[polaroid]   (12, -2,  -1.5) -- (12, -2,   1.5)  %Polarizing filter
        -- (12, 2, 1.5) -- (12, 2, -1.5) -- cycle;
    \node[anchor=north west] at (12,-2,-1.7) {Polaroid};
    \draw[ms structure, <->] (12, -1.5,-0.5) -- (12, -1.5, 0.5); %Polarization direction

    %Light leaving the second polaroid
    \draw[wave=msColor1,variable=\x, samples at={12, 12.25,..., 14}]
        plot (\x,{0}, {0.966*0.966*0.26*sin(2*(\x-0.5) r)}); %n'g polarized ray
    \draw[wave=msColor5,  variable=\x, samples at={12, 12.25,..., 14}]
        plot (\x,{0}, {-0.26*0.966*sin(2*(\x-0.1) r)});      %n'p polarized ray

    \node[ms small,align=left, text width=14cm, anchor=north west, yshift=-\msSpaceGroup] at (current bounding box.south west)
        {Light behavior in a petrographic microscope with light polarizing
        device. Only one incident wavelength is shown (monochromatic light).
        The magnetic field, perpendicular to the electric one, is not drawn.};
\end{tikzpicture}
\end{document}
